% ============================================================================
%  Chapter 11 -- Combinational Machines                          (Lecture 11)
% ============================================================================
\chapter{Combinational Machines}
\label{ch:combinational}

\epigraphline{\ldots it's not my fault, the carrots touched the peas\ldots}
  {Alan Turing as a boy, in the film \emph{The Imitation Game}, 2014}

\begin{chapgoals}
By the end of this chapter you will be able to
\begin{itemize}
\item tell a combinational machine from a sequential one;
\item compute how many bits a fixed-length code needs, and recognise incomplete and redundant codes;
\item design a decoder and an encoder, and compose small decoders into larger ones;
\item explain what a multiplexer and a demultiplexer do, and build them from gates;
\item build a comparator and a parity network, and say what a parity check can and cannot detect;
\item point to each of these machines inside a \riscv processor.
\end{itemize}
\end{chapgoals}

% ----------------------------------------------------------------------------
\section{Warm-up on the machine}

\asmfile[caption={Choose one of four colours.},label=lst:cm:palette]{code/l11/c03.S}

Listing~\ref{lst:cm:palette} is the array access of
Chapter~\ref{ch:information}, used as a selection. The same selector always
gives the same colour: the result depends only on the present inputs, and
nothing is remembered from one use to the next. In hardware this is a
\term{multiplexer}: four data inputs, two selection bits, one output. A
processor contains dozens of them, and makes the selection with gates in well
under a nanosecond, where the program takes four instructions. How are they
built?

% ----------------------------------------------------------------------------
\section{Combinational machines}

\subsection{A taxonomy of digital circuits}

\begin{definition}{Combinational and sequential circuits}
In a \term{combinational} circuit the value of the outputs at a given instant
depends only on the value of the inputs at that same instant. In a
\term{sequential} circuit the outputs depend on the present inputs and on the
past history of the inputs.
\end{definition}

A combinational circuit has no memory and no feedback loops, and is completely
described by a truth table: adders, decoders, multiplexers and the ALU are
combinational. A sequential circuit needs the concepts of \emph{time} and of
\emph{state}, and is described by a state table or a state diagram: latches,
flip-flops, registers, counters and the processor as a whole are sequential.
They are the subject of Chapter~\ref{ch:sequential}.

\subsection{Designing a combinational machine}

The design follows the method of Chapter~\ref{ch:boole2}
(Figure~\ref{fig:cm:design}): from a specification in words to a truth table,
with one row per combination of the inputs and with the don't-cares; from the
table to minimal expressions, one output at a time; from the expressions to a
network of gates, or of standard blocks. Analysis runs the other way, from
the schematic to the truth table. For more than four or five inputs nobody
draws maps: designers compose standard blocks, which is why this chapter is a
catalogue of them.

\begin{figure}[htb]
\centering
\begin{tikzpicture}[node distance=5mm]
  \node[blkacc,text width=27mm,minimum height=16mm] (a) {\textbf{Specification}\\[1pt]\footnotesize in words};
  \node[blkamb,text width=27mm,minimum height=16mm,right=of a] (b) {\textbf{Truth table}\\[1pt]\footnotesize with don't-cares};
  \node[blkvio,text width=27mm,minimum height=16mm,right=of b] (c) {\textbf{Minimisation}\\[1pt]\footnotesize maps or algebra};
  \node[blkgrn,text width=27mm,minimum height=16mm,right=of c] (d) {\textbf{Network}\\[1pt]\footnotesize gates or blocks};
  \foreach \x/\y in {a/b,b/c,c/d} \draw[->,thick] (\x) -- (\y);
\end{tikzpicture}
\caption{The synthesis of a combinational machine.}
\label{fig:cm:design}
\end{figure}

\subsection{Real gates take time}

A gate does not answer instantly: its output settles after a propagation
delay (Chapter~\ref{ch:electronics}). The delay of a network is that of its
slowest path from an input to an output, the \term{critical path}. While the
signals travel, the outputs can be temporarily wrong.

\begin{prgconcept}
A combinational machine is correct only after it has settled. Something must
say when to look at its outputs: that something is the clock, and it is where
sequential machines begin.
\end{prgconcept}

In the ripple-carry adder of Chapter~\ref{ch:boole2} the critical path
crosses all thirty-two full adders, because the last carry depends on the
first. Faster adders, such as the \emph{carry look-ahead} adder, compute the
carries with additional gates arranged in a few levels: they spend area to
shorten the critical path. The function computed is the same. The clock
period of a processor must be longer than its longest critical path.

% ----------------------------------------------------------------------------
\section{Codes}

Before building machines that treat codes, it must be clear what a code is.

\subsection{The concept of information}

A piece of information has three aspects. The \term{attribute} is the meaning
associated with it: for example, the mark of the Computer Architecture exam.
The \term{type} is the set of elements among which the value is chosen: the
marks from 18 to 30, and 30 with honours. The \term{value} is the element
actually taken: 30. To give information is to choose one element in a defined
set.

The \term{cardinality} $N$ of the type is the number of its elements, here
$N=14$; a type with more elements carries more information. The same value
can be written in many ways: 3, III, three; and A, B, C, D can stand for 30
with honours, 30, 29, 28. ``30 with honours'' is not a number: it is a
fourteenth element of the type, and how to write it is a problem of encoding.
Finally, types are \emph{simple} when they are atomic, such as integers and
characters, and \emph{structured} when they combine simple ones: a position is
a latitude and a longitude.

\subsection{Alphabets and strings}

\begin{definition}{Alphabet and string}
An \term{alphabet} is a finite, non-empty set of symbols. A \term{string}, or
word, over an alphabet is a finite sequence of its symbols. The
\term{length} $|x|$ of a string $x$ is the number of its symbols; the string
of length 0 is the empty string. The set of all the strings over an alphabet
$S$, the empty one included, is written $S^{*}$.
\end{definition}

The binary alphabet is $\{0,1\}$, the decimal one $\{0,\dots,9\}$;
\texttt{01000010} is a binary string of length 8. An infinite sequence
$010101\ldots$ is not a string, and does not belong to $\{0,1\}^{*}$.

\subsection{Codes of fixed length}

\begin{definition}{Code}
A \term{code} maps every element of a source type $T$, of cardinality $N$, to
a string of symbols of a code alphabet $E$ with $k$ symbols, the \term{code
word}. In a code of \term{fixed length} every code word has the same length
$m$.
\end{definition}

\begin{theorem}{Minimum length}
With an alphabet of $k$ symbols there are $k^{m}$ words of length $m$. A
fixed-length code for $N$ elements exists if and only if $k^{m}\ge N$, so that
the minimum length is
\[ m=\lceil \log_k N\rceil,\qquad\text{and in binary}\qquad m=\lceil\log_2 N\rceil . \]
\end{theorem}

The fourteen marks of the exam need 4 bits; the 32 registers of \riscv need
exactly 5, and those five bits are the fields \code{rd}, \code{rs1} and
\code{rs2} of an instruction. A code is \term{incomplete} when $N$ is not a
power of $k$: then $k^{m}-N$ words are left unassigned. The fourteen marks
leave $16-14=2$ words unassigned.

\subsection{An incomplete code: binary-coded decimal}

The ten decimal digits need $\lceil\log_2 10\rceil=4$ bits, and six of the
sixteen words are never used (Table~\ref{tab:cm:bcd}). In the code BCD~8421,
where each digit is written as its value in binary, the unused words are 1010
to 1111. Unused words become don't-care conditions when a circuit is designed.

\begin{table}[htb]
\centering\small
\caption{Two codes for the decimal digits.}\label{tab:cm:bcd}
\begin{tabular}{@{}ccc@{\qquad\qquad}ccc@{}}
\toprule
\textbf{Digit} & \textbf{BCD 8421} & \textbf{Excess-3} & \textbf{Digit} & \textbf{BCD 8421} & \textbf{Excess-3}\\
\midrule
0 & \texttt{0000} & \texttt{0011} & 5 & \texttt{0101} & \texttt{1000}\\
1 & \texttt{0001} & \texttt{0100} & 6 & \texttt{0110} & \texttt{1001}\\
2 & \texttt{0010} & \texttt{0101} & 7 & \texttt{0111} & \texttt{1010}\\
3 & \texttt{0011} & \texttt{0110} & 8 & \texttt{1000} & \texttt{1011}\\
4 & \texttt{0100} & \texttt{0111} & 9 & \texttt{1001} & \texttt{1100}\\
\bottomrule
\end{tabular}
\end{table}

\subsection{Redundant codes}

\begin{definition}{Redundant code}
A code is \term{redundant} when it uses more symbols than strictly necessary:
its length is greater than the minimum $m$.
\end{definition}

In a redundant code only some of the possible words are legal; the others
cannot be produced by the encoder. If a word arrives that is not legal,
something has altered it on the way: the error is \emph{detected}.
Redundancy is the price paid to detect, and sometimes to correct, errors.

\begin{center}\small
\begin{tabular}{@{}ccc@{}}
\toprule
\textbf{Element} & \textbf{Code 1} & \textbf{Code 2}\\
\midrule
A & \texttt{00} & \texttt{000}\\
B & \texttt{01} & \texttt{011}\\
C & \texttt{10} & \texttt{101}\\
D & \texttt{11} & \texttt{110}\\
\bottomrule
\end{tabular}
\end{center}

Code~1 is not redundant: every word of two bits is legal, so that an altered
bit turns one element into another without anyone noticing. Code~2 adds a
third bit that makes the number of 1s even. The words 001, 010, 100 and 111
are not legal, and a single altered bit is always detected. Code~2 is Code~1
plus a bit of \emph{even parity}; the machines that compute it come later in
the chapter.

\subsection{Codes of variable length}

In a code of \term{variable length} the length of the code word depends on the
element: frequent elements receive short words. Such a code must have the
\term{prefix property}: no code word is the beginning of a longer one, so that
the reader always knows where a word ends. It pays when the elements do not
all have the same probability: less space in memory, less time in
transmission. Morse code and Huffman coding are the classic examples, and the
reader has already met two more.

\begin{itemize}
\item \textbf{UTF-8}: one to four bytes per character
  (Chapter~\ref{ch:information}).
\item \textbf{RISC-V instructions on the kit}: with the C extension each
  instruction is 16 or 32 bits long. An instruction whose two least
  significant bits are \texttt{11} is 32 bits long; any other pattern marks an
  instruction of 16 bits. The length is written in the code word itself.
\end{itemize}

% ----------------------------------------------------------------------------
\section{Decoders and encoders}

\subsection{The decoded representation}

At the opposite extreme from the compact code there is the \term{decoded}
representation, also called \emph{one-hot}: a code in which the length of the
word equals the cardinality of the set, $m=N$, and every code word has exactly
one bit at 1. There is one line per element: \texttt{0010} means that the
third of four elements is the one chosen. The representation is simple to
use, because the active line directly selects something, but costly in wires.
The rule of use is: a compact code inside a word, a decoded representation
where a choice must be acted upon. The seven lamps of
Chapter~\ref{ch:boole1} are a decoded representation of a bit number.

\subsection{The decoder}

\begin{definition}{Decoder}
A \term{decoder} receives a code word of $m$ bits and asserts the one output
line, out of $N=2^{m}$, that corresponds to it.
\end{definition}

Each output is a minterm of the inputs: a decoder generates all of them
(Figure~\ref{fig:cm:decoder}). With an \term{enable} input EN, no output is
asserted when $\mathrm{EN}=0$.

\begin{figure}[htb]
\centering
\begin{minipage}{0.4\textwidth}\centering\small\ttfamily
\begin{tabular}{cc|cccc}
\textrm{$B$}&\textrm{$A$}&\textrm{$U_0$}&\textrm{$U_1$}&\textrm{$U_2$}&\textrm{$U_3$}\\\hline
0&0&\textbf{1}&0&0&0\\0&1&0&\textbf{1}&0&0\\1&0&0&0&\textbf{1}&0\\1&1&0&0&0&\textbf{1}
\end{tabular}
\[\begin{aligned}U_0&=\nott{B}\,\nott{A} & U_1&=\nott{B}\,A\\ U_2&=B\,\nott{A} & U_3&=B\,A\end{aligned}\]
\end{minipage}\hfill
\begin{minipage}{0.56\textwidth}\centering
\begin{circuitikz}[scale=0.8,transform shape]
  \draw (0,4.6) node[above]{$B$} -- (0,-0.4);
  \draw (1.6,4.6) node[above]{$A$} -- (1.6,-0.4);
  \draw (0,4.2) node[circ]{} -- (0.35,4.2) node[not port,anchor=in,scale=0.6](nb){};
  \draw (nb.out) -| (0.95,-0.4);
  \draw (1.6,4.2) node[circ]{} -- (1.95,4.2) node[not port,anchor=in,scale=0.6](na){};
  \draw (na.out) -| (2.55,-0.4);
  \foreach \k/\y/\xb/\xa in {0/3.2/0.95/2.55, 1/2.2/0.95/1.6, 2/1.2/0/2.55, 3/0.2/0/1.6}{
    \node[and port] (g\k) at (5,\y) {};
    \draw (\xb,0|-g\k.in 1) node[circ]{} -- (g\k.in 1);
    \draw (\xa,0|-g\k.in 2) node[circ]{} -- (g\k.in 2);
    \draw (g\k.out) -- ++(0.3,0) node[right]{$U_\k$};
  }
\end{circuitikz}
\end{minipage}
\caption{The 2-to-4 decoder ($A$ is the least significant bit).}
\label{fig:cm:decoder}
\end{figure}

\subsection{Composing decoders}

Large decoders are built from small ones. Figure~\ref{fig:cm:dec16} builds a
4-to-16 decoder from five 2-to-4 decoders with enable: the first decodes the
two high bits and enables one of the four others; the enabled one decodes the
two low bits, which go to all four, and asserts one of its four lines. Five
small decoders select one line of sixteen. Standard parts existed for this:
the SN74139 holds two 2-to-4 decoders with enable in one package. (In the
real part the enable and the outputs are active-low, a detail omitted here.)

\begin{figure}[htb]
\centering
\begin{tikzpicture}[d/.style={blkacc,minimum width=30mm,minimum height=8mm}]
  \node[blkamb,minimum width=30mm,minimum height=12mm] (h) at (0,-1.65) {decoder 2-to-4\\\footnotesize high bits $D\,C$};
  \draw[<-] (h.west) -- ++(-0.8,0) node[left,font=\ttfamily]{D C};
  \foreach \i/\a/\b in {0/0/3, 1/4/7, 2/8/11, 3/12/15}{
    \node[d] (d\i) at (5.2,-\i*1.1) {decoder 2-to-4\quad{\footnotesize EN}};
    \draw[->,prgamber,thick] ([yshift=4.5mm-\i*3mm]h.east) -- ++(0.5,0) |- (d\i.west);
    \draw[->] (d\i.east) -- ++(0.8,0) node[right,font=\ttfamily\small]{U\a\ ... U\b};
  }
  \node[lbl,text=prgcomment] at (5.2,-4.0) {\emph{the low bits $B\,A$ go to all four}};
\end{tikzpicture}
\caption{Modular composition: a 4-to-16 decoder from five 2-to-4 decoders.}
\label{fig:cm:dec16}
\end{figure}

\subsection{In the processor: which register is written?}

\begin{prgmachine}
The instruction \code{addi t1, t1, 1}, the number \hex{00130313} of
Chapter~\ref{ch:firstcontact}, contains the field $\code{rd}=\texttt{00110}$:
a code word of five bits, because there are thirty-two registers and
$\lceil\log_2 32\rceil=5$. A 5-to-32 decoder turns it into thirty-two
\emph{write-enable} lines, of which only that of \reg{x6}, that is \reg{t1},
is asserted: only that register receives the result. The write enable of
\reg{x0} is simply not connected, and \reg{x0} always reads zero. Two more
selections, on \code{rs1} and \code{rs2}, choose which registers are read:
they are made by multiplexers.
\end{prgmachine}

\subsection{The encoder}

The \term{encoder} is the inverse machine: it receives a decoded
representation, $N$ lines of which one is asserted, and outputs the
corresponding code word of $m=\lceil\log_2 N\rceil$ bits. Each output bit is
the OR of the input lines whose number has a 1 in that position; for $N=4$,
$C_0=U_1+U_3$ and $C_1=U_2+U_3$. If several lines can be asserted together, a
\term{priority encoder} is needed, which outputs the number of the asserted
line with the highest priority. A keypad with one wire per key becomes, in
this way, a compact key number; and priority encoders will appear again with
interrupts (Chapter~\ref{ch:instructions}), where several devices ask for
attention and one number reaches the processor.

\subsection{Transcoders: the seven-segment display}

A \term{transcoder} turns each word of a code $C_1$ into the corresponding
word of a code $C_2$ defined on the same set. The classic example passes from
BCD to the seven lines that light the segments $a$ to $g$ of a numeric display
(Figure~\ref{fig:cm:seg}); each segment is an LED switched by a digital
signal. The transcoder has seven outputs: seven functions of four variables
to synthesise, all with the same six don't-cares.

\begin{figure}[htb]
\centering
\begin{minipage}{0.22\textwidth}\centering
\begin{tikzpicture}[x=6mm,y=6mm,seg/.style={line width=3.2pt,line cap=round}]
  \draw[seg,prgred] (0.3,6) -- (2.7,6) node[midway,above,lbl]{a};
  \draw[seg,prgred] (3,5.7) -- (3,3.3) node[midway,right,lbl]{b};
  \draw[seg,prgrule] (3,2.7) -- (3,0.3) node[midway,right,lbl]{c};
  \draw[seg,prgred] (0.3,0) -- (2.7,0) node[midway,below,lbl]{d};
  \draw[seg,prgred] (0,2.7) -- (0,0.3) node[midway,left,lbl]{e};
  \draw[seg,prgrule] (0,5.7) -- (0,3.3) node[midway,left,lbl]{f};
  \draw[seg,prgred] (0.3,3) -- (2.7,3) node[midway,above,lbl]{g};
\end{tikzpicture}
\end{minipage}
\begin{minipage}{0.38\textwidth}\centering\scriptsize\ttfamily\setlength{\tabcolsep}{2.6pt}
\begin{tabular}{cc|ccccccc}
\textrm{digit}&\textrm{BCD}&a&b&c&d&e&f&g\\\hline
0&0000&1&1&1&1&1&1&0\\1&0001&0&1&1&0&0&0&0\\\rowcolor{prgred!12}2&0010&1&1&0&1&1&0&1\\3&0011&1&1&1&1&0&0&1\\4&0100&0&1&1&0&0&1&1\\
5&0101&1&0&1&1&0&1&1\\6&0110&1&0&1&1&1&1&1\\7&0111&1&1&1&0&0&0&0\\8&1000&1&1&1&1&1&1&1\\9&1001&1&1&1&1&0&1&1
\end{tabular}
\end{minipage}
\begin{minipage}{0.36\textwidth}\centering
\resizebox{\textwidth}{!}{\begin{karnaugh-map}[4][4][1][$CD$][$AB$]
  \minterms{0,2,3,5,6,7,8,9}\indeterminants{10,11,12,13,14,15}\autoterms[0]
  \implicant{12}{10}\implicant{3}{10}\implicant{5}{15}\implicantcorner
\end{karnaugh-map}}
\end{minipage}
\caption{The transcoder from BCD to seven segments: the digit 2 lights $a$,
$b$, $g$, $e$ and $d$. Right: the map of segment $a$, with the BCD digit on
$A\,B\,C\,D$ and $A$ most significant.}
\label{fig:cm:seg}
\end{figure}

\begin{example}{Synthesising segment $a$}
Segment $a$ is off only for the digits 1 and 4. The words 1010 to 1111 never
occur, and are marked as don't-cares in the map of
Figure~\ref{fig:cm:seg}. Using those cells as 1s the groups become large: the
two lower rows give $A$; the two right-hand columns give $C$; the central
square gives $BD$; the four corners give $\nott{B}\,\nott{D}$:
\[ a = A + C + B\,D + \nott{B}\,\nott{D}. \]
Without the don't-cares the groups would be smaller and the expression
longer. As a check, the digit 0 (\texttt{0000}) is covered by the corners and
the digit 9 (\texttt{1001}) by $A$. Six more maps give the segments $b$ to
$g$, and a complete transcoder.
\end{example}

% ----------------------------------------------------------------------------
\section{Multiplexers and demultiplexers}

\subsection{The multiplexer}

\begin{definition}{Multiplexer}
A \term{multiplexer} is a combinational circuit with $2^{n}$ data inputs, $n$
selection lines and one output. It selects one of the $2^{n}$ input signals
and transfers it to the output.
\end{definition}

Each data input is multiplied by one minterm of the selection lines, and an OR
collects the results (Figure~\ref{fig:cm:mux}):
\begin{align*}
\text{2 inputs:}\quad & Y=\nott{S}\,I_0+S\,I_1,\\
\text{4 inputs:}\quad & Y=\nott{S_1}\,\nott{S_0}\,I_0+\nott{S_1}\,S_0\,I_1+S_1\,\nott{S_0}\,I_2+S_1\,S_0\,I_3 .
\end{align*}
In other words, a decoder on the selection lines enables one AND gate: the
datum is chosen by an \emph{address}, not by a decoded line. A multiplexer
can also implement any function of its selection variables: it is enough to
put the truth table on the data inputs.

\begin{figure}[htb]
\centering
\begin{subfigure}[b]{0.48\textwidth}\centering
\begin{circuitikz}[scale=0.8,transform shape]
  \node[and port] (a0) at (3,2) {};
  \node[and port] (a1) at (3,0) {};
  \node[or port] (o) at (5.2,1) {};
  \draw (a0.in 1) -- ++(-2.2,0) node[left]{$I_0$};
  \draw (a1.in 2) -- ++(-2.2,0) node[left]{$I_1$};
  \draw (a1.in 1) -- ++(-1.6,0) coordinate (s) -- ++(0,-1.2) node[below]{$S$};
  \draw (s) node[circ]{} |- ++(0.2,1.2) node[not port,anchor=in,scale=0.6](n){};
  \draw (n.out) |- (a0.in 2);
  \draw (a0.out) -| (o.in 1); \draw (a1.out) -| (o.in 2);
  \draw (o.out) -- ++(0.3,0) node[right]{$Y$};
\end{circuitikz}
\caption{A multiplexer with two inputs}\label{fig:cm:mux}
\end{subfigure}\hfill
\begin{subfigure}[b]{0.48\textwidth}\centering
\begin{circuitikz}[scale=0.8,transform shape]
  \node[xnor port] (x0) at (2.5,2) {};
  \node[xnor port] (x1) at (2.5,0) {};
  \node[and port] (a) at (5,1) {};
  \draw (x0.in 1) node[left]{$A_0$} (x0.in 2) node[left]{$B_0$};
  \draw (x1.in 1) node[left]{$A_1$} (x1.in 2) node[left]{$B_1$};
  \draw (x0.out) -| (a.in 1); \draw (x1.out) -| (a.in 2);
  \draw (a.out) -- ++(0.3,0) node[right]{EQ};
\end{circuitikz}
\caption{A comparator for words of two bits}\label{fig:cm:cmp}
\end{subfigure}
\caption{Two machines built from gates.}
\end{figure}

\subsection{The demultiplexer}

The \term{demultiplexer} is the inverse machine: one data input $D$, $n$
selection lines, $2^{n}$ outputs. It routes the input to the output chosen by
the selection lines, and the other outputs stay at 0:
\[ F_i = D\cdot U_i , \]
where $U_i$ are the outputs of a decoder on the selection lines. With $D$
fixed at 1 a demultiplexer is exactly a decoder, in which $D$ plays the part
of the enable. A multiplexer followed by a demultiplexer carries many signals
over one wire, one at a time: the idea behind every shared line.

\subsection{In the processor: register or constant?}

\begin{prgmachine}
The second operand of the ALU is either the value of the register \code{rs2}
or the constant written in the instruction, sign-extended
(Figure~\ref{fig:cm:alusrc}). A multiplexer of two inputs, thirty-two bits
wide, chooses: its selection line is 0 for \ins{add} and 1 for \ins{addi}. In
the datapath of \citet{patterson2020} the signal is called \code{ALUSrc}. A
gate applied to every bit of a word under a common control is called a
\emph{word gate}.
\end{prgmachine}

\begin{figure}[htb]
\centering
\begin{tikzpicture}
  \node[blk,fill=prgsteel,text=white,minimum width=34mm] (r1) at (0,2.2) {register file: \texttt{rs1}};
  \node[blk,fill=prgsteel,text=white,minimum width=34mm] (r2) at (0,0.9) {register file: \texttt{rs2}};
  \node[blkamb,minimum width=34mm] (im) at (0,-0.3) {immediate, sign-extended};
  \node[blkacc,minimum width=10mm,minimum height=20mm] (m) at (3.6,0.3) {MUX};
  \node[blkvio,minimum width=18mm,minimum height=30mm] (alu) at (6.4,1.2) {\textbf{ALU}};
  \draw[->,bus] (r1.east) -- (r1.east-|alu.west) node[pos=0.9,above,lbl]{A};
  \draw[->,bus] (r2.east) -- (r2.east-|m.west) node[pos=0.85,above,lbl]{0};
  \draw[->,bus] (im.east) -- (im.east-|m.west) node[pos=0.85,above,lbl]{1};
  \draw[->,bus] (m.east) -- (m.east-|alu.west) node[pos=0.8,above,lbl]{B};
  \draw[->,bus] (alu.east) -- ++(1,0) node[right,lbl]{result};
  \draw[->,thick,prgred] (3.6,-1.5) node[below,lbl,text=prgred]{selection: 0 for \texttt{add}, 1 for \texttt{addi}} -- (m.south);
\end{tikzpicture}
\caption{A multiplexer chooses the second operand of the ALU.}
\label{fig:cm:alusrc}
\end{figure}

% ----------------------------------------------------------------------------
\section{Comparing and checking}

\subsection{The comparator}

A \term{comparator} has two data $A$ and $B$ in input and one Boolean output
EQ that is 1 when $A=B$. Two bits are equal when their XNOR is 1; two words
are equal when all their pairs of bits are equal, an AND of XNORs
(Figure~\ref{fig:cm:cmp}):
\[ \mathrm{EQ}=\nott{(A_0\oplus B_0)}\cdot\nott{(A_1\oplus B_1)}\cdots \]
In the processor this is \ins{beq} and \ins{bne}: the two registers are
compared bit by bit, and the branch is taken or not according to EQ. Many
datapaths reuse the ALU instead: they subtract and test whether the result is
zero. Comparisons of order (\ins{blt}, \ins{slt}) need the subtraction and
the sign rules of Chapter~\ref{ch:integers}.

\subsection{Parity}

When data travel from a source to a destination, redundant codes allow errors
to be checked. The simplest check is a single \term{parity bit}: the legal
code words have, for example, an even number of bits at 1.

\begin{definition}{Parity network}
A \term{parity network} computes the parity of the $n$ bits at its input,
\[ p=x_0\oplus x_1\oplus\dots\oplus x_{n-1}, \]
which is 1 when the number of 1s is odd.
\end{definition}

XOR is addition modulo 2, which is why a chain of XORs counts the 1s modulo 2.
The network is needed twice (Figure~\ref{fig:cm:parity}): at the source, to
compute the bit to be added to a non-redundant code, and at the destination,
to verify the overall parity of the redundant code on arrival.

\begin{figure}[htb]
\centering
\begin{tikzpicture}[node distance=4mm]
  \node[blkacc,text width=21mm,minimum height=16mm] (a) {\textbf{Source}\\[1pt]\footnotesize code $C$};
  \node[blkamb,text width=21mm,minimum height=16mm,right=of a] (b) {\textbf{CALC}\\[1pt]\footnotesize appends the parity bit};
  \node[blkvio,text width=21mm,minimum height=16mm,right=of b] (c) {\textbf{Channel}\\[1pt]\footnotesize may alter bits};
  \node[blkgrn,text width=21mm,minimum height=16mm,right=of c] (d) {\textbf{VER}\\[1pt]\footnotesize verifies the parity};
  \node[blkred,text width=21mm,minimum height=16mm,right=of d] (e) {\textbf{Destination}\\[1pt]\footnotesize code $C'$};
  \foreach \x/\y in {a/b,b/c,c/d,d/e} \draw[->,thick] (\x) -- (\y);
  \draw[->,thick,prgred] (d.south) -- ++(0,-5mm) node[below,lbl,text=prgred]{error signal};
\end{tikzpicture}
\caption{Transmission with a parity check.}
\label{fig:cm:parity}
\end{figure}

\begin{theorem}{Power of a parity bit}
A single parity bit detects every alteration of an odd number of bits. It
does not detect the alteration of an even number of bits, and it cannot
correct any error.
\end{theorem}

\noindent\emph{Proof.} Each altered bit changes the number of 1s by one, and
therefore inverts the parity: an odd number of inversions changes it, an even
number leaves it as it was. And the check gives one bit of information, right
or wrong, which is not enough to say \emph{which} bit was altered.
\hfill$\square$

\smallskip
The code $C'$ that reaches the destination is thus \emph{probably} equal to
$C$ if no error was signalled. Stronger codes add more check bits: the codes
of \citet{hamming1950}, with several parity bits over overlapping groups of
bits, locate and correct a single error.

\subsection{A parity network in seven instructions}

\asmfile[caption={The parity of a byte.},label=lst:cm:parity]{code/l11/c06.S}

Each \ins{xor} of Listing~\ref{lst:cm:parity} combines the upper half of the
bits with the lower half: three levels for eight bits, the software version of
a tree of seven XOR gates. The value computed, appended to the byte, makes
its total parity even. Trace it with $\hex{53}=0101\,0011_2$, which has four
1s and gives 0, and with $\hex{13}=0001\,0011_2$, which has three and gives 1.

% ----------------------------------------------------------------------------
\section{The arithmetic logic unit}

All the blocks of this chapter work together in one machine.

\begin{definition}{Arithmetic logic unit}
The \term{arithmetic logic unit} (ALU) is a combinational machine that
receives two operands and a code of operation, and produces the result of the
operation selected.
\end{definition}

\begin{figure}[htb]
\centering
\begin{tikzpicture}[u/.style={blk,minimum width=40mm,minimum height=7mm}]
  \foreach \i/\t/\c in {0/{adder / subtractor}/prgviolet!12, 1/{AND, one gate per bit}/prgpanel, 2/{OR, one gate per bit}/prgpanel, 3/{XOR, one gate per bit}/prgpanel, 4/{shifter}/prgpanel}{
    \node[u,fill=\c] (u\i) at (0,-\i*0.95) {\t};
    \draw[->,prgaccent,thick] (-3.6,-\i*0.95+0.12) -- ([yshift=1.2mm]u\i.west);
    \draw[->,prgamber,thick] (-3.0,-\i*0.95-0.12) -- ([yshift=-1.2mm]u\i.west);
    \draw[->] (u\i.east) -- (3.4,-\i*0.95);
  }
  \draw[prgaccent,very thick] (-3.6,0.6) node[above,font=\ttfamily\bfseries]{A} -- (-3.6,-3.68);
  \draw[prgamber,very thick] (-3.0,0.6) node[above,font=\ttfamily\bfseries]{B} -- (-3.0,-3.92);
  \node[blkacc,minimum width=11mm,minimum height=46mm] (m) at (3.95,-1.9) {MUX};
  \draw[->,bus] (m.east) -- ++(1.1,0) node[right,lbl]{result};
  \draw[->,thick,prgred] (3.95,-5.0) node[below,lbl,text=prgred]{operation, from \texttt{funct3}} -- (m.south);
\end{tikzpicture}
\caption{The ALU is a set of machines, all at work on the same two operands of
32 bits, and a multiplexer that chooses one result.}
\label{fig:cm:alu}
\end{figure}

The structure of Figure~\ref{fig:cm:alu} is disarmingly simple: \emph{every}
unit computes its result at every instruction, and a multiplexer lets one
through. The adder is the ripple-carry adder of Chapter~\ref{ch:boole2}, with
the trick of the complement and of the carry-in of
Chapter~\ref{ch:integers} for the subtraction; comparisons (\ins{slt}) reuse
the subtractor.

The selection is written in the instruction (Table~\ref{tab:cm:funct3}): the
field \code{funct3}, bits 14 to 12, is a compact code of three bits, eight
code words, which a small decoder turns into the selection lines of the
multiplexer of the ALU.

\begin{table}[htb]
\centering\small
\caption{\riscv instructions and the ALU.}\label{tab:cm:funct3}
\begin{tabularx}{\textwidth}{@{}lllL@{}}
\toprule
\textbf{Instruction} & \textbf{funct3} & \textbf{ALU operation} & \textbf{Unit that produces the result}\\
\midrule
\ins{add}, \ins{sub} & \texttt{000} & sum or difference & adder / subtractor (one more bit, in \code{funct7}, chooses)\\
\ins{sll} & \texttt{001} & shift left & shifter\\
\ins{slt}, \ins{sltu} & \texttt{010}, \texttt{011} & set if less than & subtractor, then the sign rules\\
\ins{xor} & \texttt{100} & bitwise XOR & XOR gates\\
\ins{srl}, \ins{sra} & \texttt{101} & shift right & shifter (\code{funct7} chooses logical or arithmetic)\\
\ins{or} & \texttt{110} & bitwise OR & OR gates\\
\ins{and} & \texttt{111} & bitwise AND & AND gates\\
\bottomrule
\end{tabularx}
\end{table}

% ----------------------------------------------------------------------------
% Generated by tools/textbook/export_listings.js from decks/plan.js and code/l11.S: do not edit.
\section{The program of this chapter}\label{sec:lab:l11}

The listings of this chapter are parts of one program, \fname{code/l11.S} in the
repository of the course, printed in full in Listing~\ref{lst:lab:l11}. The
automatic tests of the course material build it and run it on the headless
runner of \prgqt{}, with scripted input, and check the frame;
Figure~\ref{fig:lab:l11} is the frame that the test wrote.

\begin{description}[style=nextline,leftmargin=1.5em]
\item[On the screen.] Four horizontal bars, red, green, blue and yellow from the top: bar k is drawn in the colour that the selector k chooses in the table.
\item[What to press.] Nothing: the program runs by itself.
\end{description}

\begin{figure}[htb]
\centering
\includegraphics[width=0.5\textwidth]{lab_l11.png}
\caption{\fname{l11.S} on \prgqt{}.}\label{fig:lab:l11}
\end{figure}

\asmfile[caption={The complete program of the chapter, \fname{code/l11.S}.},label=lst:lab:l11]{code/l11/program.S}

To run it on \prgqt{} the commands are those of Listing~\ref{lst:lab:l11:run};
Section~\ref{sec:prg:testing} gives the same steps for QEMU and for the board.

\begin{shell}[caption={Building \fname{l11.S} and running it on \prgqt{}.},label=lst:lab:l11:run]
# 1. start PRG32-QT on this computer (it listens on port 8080)

# 2. in the PRG32 folder, with ESP-IDF active: build the cartridge
python3 -m prg32 cartridge build ~/arc12/code/l11.S --portable \
    --entry-prefix l11 --name l11 --out l11.prg32

# 3. send it to PRG32-QT and start it
python3 -m prg32 esp32c6 upload-and-run l11.prg32 \
    --url http://localhost:8080 --slot cart0

# 4. look, press, compare with the previous slide
\end{shell}


% ----------------------------------------------------------------------------
\section{Check your understanding}

\begin{quiz}{How many bits does a fixed-length binary code need for the 26 letters of the alphabet?}
\item 4
\item 5
\item 6
\item 26
\end{quiz}

\begin{quiz}{A decoder has 3 input lines and no enable. How many outputs does it have, and how many are 1 at any time?}
\item 3 outputs, one at 1
\item 6 outputs, two at 1
\item 8 outputs, one at 1
\item 8 outputs, three at 1
\end{quiz}

\begin{quiz}{A multiplexer with 8 data inputs needs how many selection lines?}
\item 2
\item 3
\item 4
\item 8
\end{quiz}

\begin{quiz}{A word with one even-parity bit is received with two bits altered. The parity check:}
\item detects the error and corrects it
\item detects the error but cannot correct it
\item does not detect the error
\item signals an error only if both bits were 1
\end{quiz}

\section{Exercises}

\begin{exercise}[paper first, then \prg]
\begin{enumerate}[nosep]
\item Synthesise segment $g$ of the seven-segment transcoder with a Karnaugh
  map, using the don't-cares.
\item Build a 3-to-8 decoder from two 2-to-4 decoders with enable and one
  inverter.
\item Implement $F=\sum m(1,2,4,7)$ with one multiplexer of 8 inputs. Which
  function is it?
\item Write a comparator in assembly: $\reg{a0}=1$ if \reg{a0} and \reg{a1}
  are equal, using \ins{xor} and \ins{sltiu} only.
\item Draw a seven-segment digit on \prg: a table of ten bytes, one per digit,
  and seven rectangles lit by testing its bits.
\end{enumerate}
\expected{The third task puts 0, 1, 1, 0, 1, 0, 0, 1 on the data inputs: it is
the XOR of three inputs, the sum bit of the full adder. The fourth is
\code{xor a0, a0, a1} followed by \code{sltiu a0, a0, 1} (the
pseudo-instruction \ins{seqz}). On \prg: a large digit that counts from 0 to
9, one step per second, drawn from seven rectangles; it is a software
transcoder followed by seven bit tests, and its table is the truth table of
Figure~\ref{fig:cm:seg}, one row per byte.}
\end{exercise}

\section{Glossary}

\begin{glossarytable}
combinational network & rete combinatoria\\
propagation delay & ritardo di propagazione\\
cardinality & cardinalit\`a\\
code word & parola codice\\
fixed / variable length & lunghezza fissa / variabile\\
redundant code & codice ridondante\\
decoder / encoder & decodificatore / codificatore\\
transcoder & transcodificatore\\
multiplexer / demultiplexer & multiplexer / demultiplexer\\
enable & abilitazione\\
comparator & comparatore\\
parity bit & bit di parit\`a\\
\end{glossarytable}

\section{Conclusions and reading}

Information is an attribute, a type and a value; a code writes the value as a
string over an alphabet. Fixed-length codes need $\lceil\log_2 N\rceil$ bits;
redundant codes spend more to detect errors. Encoding, decoding and
transcoding are done by combinational machines. Multiplexers select a source,
demultiplexers a destination; comparators and parity networks check. A
\riscv processor uses all of them: decoders on the fields of the instruction,
multiplexers on the operands, and an ALU that is itself built from these
blocks.

What is missing here is memory. Before the sequential machines of
Chapter~\ref{ch:sequential}, the next chapter puts the logic of this part to
work in the second part of \vermone.

\paragraph{Reading.} \citet{patterson2020}, appendix~A and section~4.1;
\citet{bucci2017}, the chapters on combinational networks;
\citet{stallings2021}, the chapter on digital logic.
